\begin{table}[t]
\centering
\caption{Species-Level Mechanism Test: Targeted vs.\ Non-Targeted Fish}
\label{tab:ddd}
\begin{threeparttable}
\begin{tabular}{lc}
\toprule
 & Log Fish Density \\
\midrule
MPA $\times$ Post & NA \\
 & (NA) \\
\addlinespace
MPA $\times$ Post $\times$ Targeted & 0.059** \\
 & (0.018) \\
\addlinespace
Site $\times$ Species FE & Yes \\
Year $\times$ Species FE & Yes \\
Site $\times$ Year FE & Yes \\
Observations & 15,075 \\
\bottomrule
\end{tabular}
\begin{tablenotes}[flushleft]
\footnotesize
\item \textit{Notes:} Triple-difference specification at the site $\times$ year $\times$ species level. The dependent variable is log fish density (count per 60m$^2$ + 1). ``Targeted'' indicates species commonly taken by recreational or commercial fishing in California kelp forests (genera \textit{Semicossyphus}, \textit{Paralabrax}, \textit{Sebastes}, \textit{Scorpaena}, \textit{Caulolatilus}, \textit{Rhacochilus}, plus surfperch, cabezon, and lingcod). The triple interaction captures whether targeted species differentially recover in MPA sites after designation, controlling for site-specific species baselines, common species trends, and site-level shocks. Standard errors clustered at the site level.
\end{tablenotes}
\end{threeparttable}
\end{table}
